The United States Census Bureau counts incarcerated people at the location of
their prison or jail rather than their pre-incarceration home address---a
practice known as prison gerrymandering. Because approximately 58\% of the
state and federal prison population is Black or Hispanic, and because large
correctional facilities are concentrated in rural areas that tend to vote
Republican, this Census methodology systematically inflates the population
count of rural legislative districts while deflating the counts of the urban
and suburban communities from which most incarcerated people originate.
At least twelve states have enacted legislation to adjust prison populations
to home addresses for redistricting purposes, using data now available through
the Census Bureau's 2020 Detailed Demographic and Housing Characteristics
(DHC-A) file.

This paper analyzes the legal framework for prison gerrymandering, surveys
enacted state adjustment laws, and documents how the \textsc{bisect}
algorithmic redistricting pipeline handles prison population adjustments.
The \textsc{bisect-data} preprocessing module accepts adjusted Census data
where state law requires it, applying the adjustment uniformly before any
partisan data enters the pipeline. Empirical results for North Carolina,
Texas, and California show that prison population adjustment shifts district
boundaries in 1--2 districts per state, with partisan lean changes of
0.3--0.5 percentage points and Voting Rights Act
implications for minority-opportunity districts. Because \textsc{bisect}
applies adjustment uniformly and reproducibly---with a full SHA audit
chain---it provides expert witnesses and courts with a verifiable,
neutral implementation of whichever population base state law requires.
